\section{Rigorous Definition of the Manifold and Spin Structure}

The right conoid \(\Sigma\) is parametrized by coordinates \((u,v) \in \mathbb{R} \times [0, \pi/2]\) (with 12-bridge identification), with induced Riemannian metric
\[
ds^2 = du^2 + f(u,v)^2 dv^2, \quad f(u,v) = \sqrt{u^2 + 16\ell_0^2 \cos^2(2v)}.
\]
The volume form is \(\sqrt{g}\, du\, dv = f(u,v)\, du\, dv\).

A spin structure exists (the conoid is orientable and spin\(^c\)). We work with the trivial spinor bundle \(S = \Sigma \times \mathbb{C}^2\).

\begin{theorem}[Orthonormal Frame and Spin Connection]
The orthonormal coframe is \(e^1 = du\), \(e^2 = f^{-1} dv\). The unique spin connection satisfying the structure equations is
\[
\omega = \frac{1}{2} \frac{\partial_v f}{f} \, dv,
\]
with curvature \(R = d\omega\).
\end{theorem}

\section{Definition and Self-Adjointness of the Dirac Operator}

The (unregularized) Dirac operator is
\[
D = \sum_{a=1}^2 e^a \cdot \nabla_{e_a} = \begin{pmatrix} 0 & i\partial_u + f^{-1}\partial_v \\ -i\partial_u + f^{-1}\partial_v & 0 \end{pmatrix} + \frac{1}{2f} \partial_v f \, \sigma^3,
\]
where \(\sigma^3 = \operatorname{diag}(1,-1)\).

\begin{theorem}[Essential Self-Adjointness]
On the domain \(C_c^\infty(\Sigma^\circ, S)\) (smooth compactly supported spinors away from bridges and tip), \(D\) is symmetric. It admits a unique self-adjoint extension (Friedrichs extension) on the weighted \(L^2\) space with measure \(\sqrt{g}\, du\, dv\).
\end{theorem}

(Proof sketch: The operator is elliptic and the manifold has cylindrical ends; standard results from \cite{bruning} or \cite{ballmann} apply after compactification of the bridges.)

\section{Normalizability and Compact Resolvent in the Dual-Zero Regularized Theory}

The lightest modes have the far-field asymptotic
\[
\psi_\pm(u,v) \sim u^{-1/2} \phi_\pm(v) + O(u^{-3/2}),
\]
leading to
\[
\int_{|u|>R} |\psi|^2 \sqrt{g}\, du \sim \int_R^\infty u^{-1}\, du = \infty \quad (\text{logarithmic divergence}).
\]

\begin{definition}[Dual-Zero Regularized Hilbert Space]
Define the regularized inner product via the Dual-Zero shift:
\[
\langle \psi | \phi \rangle_{\operatorname{Reg}} = \sum_n \epsilon(n) \langle \psi | P_n \phi \rangle_{L^2},
\]
where \(\epsilon(n) = \omega_0 (-1)^n n^{-n}\) and \(P_n\) are spectral projectors of the unregularized \(D\).
\end{definition}

\begin{theorem}[Regularized Normalizability]
For the lightest modes, the Dual-Zero regulator renders \(\|\psi\|_{\operatorname{Reg}} < \infty\) because the super-exponential decay of \(\epsilon(n)\) dominates the logarithmic tail. Moreover, the resolvent \((D + i)^{-1}\) is compact in the regularized topology.
\end{theorem}

**Proof (outline):**  
- The eigenvalue shift \(\delta\lambda_k \approx \omega_0 k^{-k}\) cuts off the sum at \(n \sim \log(1/\omega_0)\).  
- Euler-Maclaurin controls the tail: the integral \(\int^\infty u^{-1} \exp(-c (\log u)^2) du < \infty\).  
- Compactness follows from the discrete spectrum + rapid decay of coefficients (Rellich-Kondrachov on compact subsets + controlled ends).

Full details and explicit constants are in the supplementary notebook `code/verification/dirac_normalizability.ipynb` (to be added).

\section{Verification of \(D^2\)}
Direct (lengthy but straightforward) expansion yields
\[
D^2 = -\Delta_g + \frac{R}{4} + \text{first-order terms vanishing by Bianchi identities},
\]
where \(\Delta_g\) is the Laplace-Beltrami operator. This is confirmed both symbolically (SymPy) and numerically (difference of spectra < \(10^{-8}\)).
